% Part: normal-modal-logic
% Chapter: frame-definability
% Section: standard-translation

\documentclass[../../../include/open-logic-section]{subfiles}

\begin{document}

\olfileid{nml}{frd}{st}

\olsection{Definabilitas Orde Kapindho}

Ora saben sipat bingkai kang bisa didéfinisikake dening formula
modal bisa didéfinisikake kanthi orde kapisan. Nanging, yen kita
ngidinake kuantifikasi tumrap predikat siji panggonan (yaiku
kuantifikasi orde kapindho monadik), kita bisa ndéfinisikake kabeh
sipat bingkai kang bisa didéfinisikake kanthi modal. Carane yaiku
nggunakake hubungan kang sistematis antarane kondisi kang ndadekake
!!{formula} modal bener ing sawijining donya lan !!{formula}s orde
kapisan. Iki diarani translasi standar saka formula modal menyang
formula orde kapisan ing basa kang ora mung ngemot !!{predicate} rong
panggonan~$Q$ kanggo relasi aksesibilitas, nanging uga !!{predicate}
siji panggonan~$P_i$ kanggo !!{propositional variable}s~$\Obj p_i$
kang katon ing~$!A$.

\begin{defn}
  \emph{Translasi standar}~$\ST_x(!A)$ didéfinisikake kanthi
  induksi kaya mangkene:
  \begin{enumerate}
  \tagitem{prvFalse}{\indcase{!A}{\lfalse}{$\ST_x(\indfrm) = \lfalse$.}}{}
  \tagitem{prvTrue}{\indcase{!A}{\ltrue}{$\ST_x(\indfrm) = \ltrue$.}}{}
  \item \indcase{!A}{\Obj p_i}{$\ST_x(\indfrm) = \Atom{P_i}{x}$.}
  \tagitem{prvNot}{\indcase{!A}{\lnot !B}{$\ST_x(\indfrm) = \lnot \ST_x(!B)$.}}{}
  \tagitem{prvAnd}{\indcase{!A}{(!B \land !C)}{$\ST_x(\indfrm) = (\ST_x(!B)
      \land \ST_x(!C))$.}}{}
  \tagitem{prvOr}{\indcase{!A}{(!B \lor !C)}{$\ST_x(\indfrm) = (\ST_x(!B)
      \lor \ST_x(!C))$.}}{}
  \tagitem{prvIf}{\indcase{!A}{(!B \lif !C)}{$\ST_x(\indfrm) = (\ST_x(!B)
      \lif \ST_x(!C))$.}}{}
  \tagitem{prvIff}{\indcase{!A}{(!B \liff !C)}{$\ST_x(\indfrm) = (\ST_x(!B)
      \liff \ST_x(!C))$.}}{}
  \tagitem{prvBox}{\indcase{!A}{\Box !B}{$\ST_x(\indfrm) =
      \lforall[y][(\Atom{Q}{x,y} \lif \ST_y(!B))]$.}}{}
  \tagitem{prvDiamond}{\indcase{!A}{\Diamond !B}{$\ST_x(\indfrm) =
      \lexists[y][(\Atom{Q}{x,y} \land \ST_y(!B))]$.}}{}
  \end{enumerate}
\end{defn}

Upamane, $\ST_x(\Box p \lif p)$ yaiku $\lforall[y][(\Atom{Q}{x,y}
  \lif \Atom{P}{y})] \lif \Atom{P}{x}$. Saben !!{structure} kanggo
basa~$\ST_x(!A)$ mbutuhake domain, relasi rong panggonan kang
ditetepake kanggo~$Q$, lan subhimpunan domain kang ditetepake kanggo
!!{predicate}s siji panggonan~$P_i$. Kanthi tembung liya, komponen
sawijining !!a{structure} kaya mangkono persis komponen model kanggo~$!A$:
domain iku himpunan donya, relasi rong panggonan kang ditetepake
kanggo~$Q$ iku relasi aksesibilitas, lan subhimpunan kang ditetepake
kanggo~$P_i$ iku penetapan~$V(\Obj p_i)$. Mula ora nggumunake yen
kepuasan $!A$ ing model modal lan kepuasan $\ST_x(!A)$ ing
!!{structure} kang cocog iku padha:

\begin{prop}\ollabel{prop:st}
  Upamane $\mModel{M} = \tuple{W, R, V}$, lan $\Struct{M'}$
  sawijining !!{structure} orde kapisan kanthi $\Domain{M'} = W$,
  $\Assign{Q}{M'} = R$, lan $\Assign{P_i}{M'} = V(\Obj p_i)$,
  sarta $s(x) = w$. Mula
  \[
  \mSat{M}{!A}[w] \text{ yen lan mung yen } \Sat{M'}{\ST_x(!A)}[s]
  \]
\end{prop}

\begin{proof}
  Kanthi induksi marang~$!A$.
\end{proof}

\begin{prop}
  Upamane $!A$ iku !!{formula} modal lan $\mModel{F} = \tuple{W, R}$
  sawijining bingkai. Upamane $\Struct{F'}$ sawijining !!{structure}
  orde kapisan kanthi $\Domain{F'} = W$ lan $\Assign{Q}{F'} = R$,
  lan $!A'$ formula orde kapindho iki:
  \[
  \lforall[X_1][\dots\lforall[X_n][\lforall[x][
        \SSubst{\ST_x(!A)}{\subst{X_1}{P_1}, \dots,
          \subst{X_n}{P_n}}]]],
  \]
  ing ngendi $P_1$, \dots, $P_n$ kabeh !!{predicate}s siji panggonan
  ing~$\ST_x(!A)$. Mula
  \[
  \mModel{F} \Entails !A \text{ yen lan mung yen } \Sat{F'}{!A'}
  \]
\end{prop}

\begin{proof}
  $\Sat{F'}{!A'}$ yen lan mung yen kanggo saben !!{structure}~$\mModel{M'}$
  kang $\Assign{P_i}{M'} \subseteq W$ kanggo $i = 1$, \dots,~$n$,
  lan kanggo saben $s$ kang $s(x) \in W$,
  $\Sat{M'}{\ST_x(!A)}[s]$. Miturut \olref{prop:st}, iki bener yen
  lan mung yen kanggo kabeh model~$\mModel{M}$ kang
  adhedhasar~$\mModel{F}$ lan saben donya~$w \in W$,
  $\mSat{M}{!A}[w]$, yaiku $\mModel{F} \Entails !A$.
\end{proof}

\begin{defn}
  Kelas bingkai~$\mClass{F}$ \emph{bisa didéfinisikake kanthi orde
    kapindho} yen ana !!a{sentence}~$!A$ ing basa orde kapindho kang
  mung nduweni siji !!{predicate} rong panggonan~$P$ lan kuantor
  mung tumrap variabel himpunan monadik, saengga $\mModel{F} =
  \tuple{W, R} \in \mClass{F}$ yen lan mung yen $\Sat{M}{!A}$ ing
  !!{structure}~$\Struct{M}$ kanthi $\Domain{M} = W$ lan
  $\Assign{P}{M} = R$.
\end{defn}

\begin{cor}
  Yen kelas bingkai bisa didéfinisikake dening !!a{formula}~$!A$,
  kelas relasi aksesibilitas kang cocog bisa didéfinisikake dening
  ukara orde kapindho monadik.
\end{cor}

\begin{proof}
  Ukara orde kapindho monadik~$!A'$ saka bukti sadurunge nduweni
  sipat kang dibutuhake.
\end{proof}

Minangka tuladha, delengen maneh !!{formula}~$\Box p \lif p$.
Formula iki ndéfinisikake refleksivitas. Refleksivitas mesthi bisa
didéfinisikake kanthi orde kapisan dening !!{sentence}~$\lforall[x][\Atom{Q}{x,x}]$.
Nanging sipat iku uga bisa didéfinisikake dening !!{sentence} orde
kapindho monadik
\[
\lforall[X][\lforall[x][(\lforall[y][(\Atom{Q}{x,y} \lif \Atom{X}{y})]
    \lif \Atom{X}{x})]].
\]
Iki mesthi ateges yen loro ukara kasebut ekuivalen. Mangkene carane
ndeleng kanthi langsung: Kaping pisan, upamane ukara orde kapindho
iku bener ing !!a{structure}~$\Struct{M}$. Amarga $x$ lan $X$
dikuwantifikasi universal, bagean kang isih ana kudu bener kanggo
sembarang $x \in W$ lan himpunan~$X \subseteq W$, upamane himpunan
$\Setabs{z}{Rxz}$ kang $R = \Assign{Q}{M}$. Dadi, kanggo saben~$s$
kang $s(x) \in W$ lan $s(X) = \Setabs{z}{Rxz}$, kita duwe
$\Sat{M}{\lforall[y][(\Atom{Q}{x,y} \lif
    \Atom{X}{y})] \lif \Atom{X}{x}}$. Nanging amarga pilihan
$s(X)$, iki ateges $\Sat{M}{\lforall[y][(\Atom{Q}{x,y} \lif
    \Atom{Q}{x,y})] \lif \Atom{Q}{x,x}}[s]$, kang ekuivalen karo
$\Atom{Q}{x,x}$ amarga antesedene valid. Amarga $s(x)$ dipilih
sebarang, kita duwe $\Sat{M}{\lforall[x][\Atom{Q}{x,x}]}$.

Saiki upamane $\Sat{M}{\lforall[x][\Atom{Q}{x,x}]}$, lan
nuduhna yen
$\Sat{M}{\lforall[X][\lforall[x][(\lforall[y][(\Atom{Q}{x,y} \lif
        \Atom{X}{y})] \lif \Atom{X}{x})]]}$. Pilih sembarang
penetapan~$s$, lan upamane $\Sat{M}{\lforall[y][(\Atom{Q}{x,y} \lif
    \Atom{X}{y})]}[s]$. Pilih $s'$ minangka varian-$y$ saka~$s$
kanthi $s'(y) = s(x)$; kita duwe
$\Sat{M}{\Atom{Q}{x,y} \lif \Atom{X}{y}}[s']$, yaiku
$\Sat{M}{\Atom{Q}{x,x} \lif \Atom{X}{x}}[s]$. Amarga
$\Sat{M}{\lforall[x][\Atom{Q}{x,x}]}$, antesedene bener, mula
$\Sat{M}{\Atom{X}{x}}[s]$, kaya kang kudu dituduhake.

Amarga saweneh kelas bingkai kang bisa didéfinisikake
ora bisa didéfinisikake kanthi orde kapisan, ora saben
!!{sentence} orde kapindho monadik kang wujude~$!A'$ ekuivalen
karo !!{sentence} orde kapisan. Ora ana cara efektif kanggo mutusake
endi kang ekuivalen.

\end{document}
